\section{同一骨架的三种模型家族}

前一章实现了 decoder-only GPT。本章只改变 connectivity 与 objective，就得到 encoder-only 和 encoder--decoder 家族。这个统一视角比死记模型名更可靠：先问每个位置能看谁，再问 keys/values 来自哪里，最后问训练目标要求输出什么。

\subsection{Encoder：取消 causal mask}

Encoder self-attention 允许每个 token 同时读取左、右上下文，因此适合生成 contextual representation、分类或 denoising。代码层面可以把 causal mask 删除，让 attention matrix 的所有位置参与 softmax；目标则不能再是同一个 left-to-right next-token loss，否则未来 token 会泄漏答案。

\lecturefigure{slide-44-encoder-blocks.jpg}{Encoder blocks：删除 causal masking，让所有 token 双向通信。}{Stanford Online 官方视频 00:43:18--00:43:53。}

\begin{knowledgebox}{读图：同一 attention，约束不同}
GPT 的下三角 connectivity 强制因果生成；BERT 类 encoder 的全连接 self-attention 让每个 token 汇聚整句信息。BERT 常用 masked/denoising objectives：先破坏输入，再要求模型恢复或提取表示。架构差别不只是“有没有 mask”，objective 也必须与可见信息匹配。
\end{knowledgebox}

\subsection{Cross-attention：decoder 查询 encoder memory}

Encoder--decoder 模型先把源序列编码成 $H_{\text{enc}}$，decoder state $X_{\text{dec}}$ 产生 queries，而 keys/values 来自 encoder：

\[
Q=X_{\text{dec}}W_Q,\qquad
K=H_{\text{enc}}W_K,\qquad
V=H_{\text{enc}}W_V.
\]

因此目标 token 既能通过 causal self-attention 看已生成前缀，也能通过 cross-attention 读取完整源输入。机器翻译、摘要和 text-to-text tasks 都可使用这种显式条件通道。

\lecturefigure{slide-45-cross-attention.jpg}{Cross-attention：decoder queries 从 encoder keys/values 读取源信息。}{Stanford Online 官方视频 00:43:53--00:44:47。}

\begin{knowledgebox}{读图：箭头方向对应信息来源}
Decoder self-attention 的 Q/K/V 都来自 decoder states；cross-attention 的 query 仍来自当前 decoder node，但 key/value 来自 encoder 顶层。两者共享 dot product、softmax 与 value aggregation。所谓“cross”不是另一种神秘算法，而是跨两组 states 进行同一寻址操作。
\end{knowledgebox}

\subsection{GPT、BERT 与 T5 的结构化对照}

模型名很多时，应把它们还原成 connectivity 与 objective。下表只描述经典原型，现实模型可能加入 span corruption、prefix LM、mixture-of-experts 或 multimodal adapters；判断新模型时仍可沿相同三问分析。

\lecturefigure{slide-46-transformer-model-families.jpg}{三大家族：decoder-only GPT、encoder-only BERT、encoder--decoder T5。}{Stanford Online 官方视频 00:44:47--00:45:31。}

\begin{table}[H]
\centering
\small
\begin{tabularx}{\textwidth}{L{2.8cm}L{3.5cm}L{3.2cm}X}
\toprule
家族 & Connectivity & 典型目标 & 适合的问题 \\
\midrule
Decoder-only / GPT & 左到右的因果自注意力 & 预测下一个 token & 开放式生成、按 prompt 补全 \\
Encoder-only / BERT & 双向自注意力 & 掩码恢复、去噪或监督任务 & 表征提取、分类、检索、序列标注 \\
Encoder--decoder / T5 & 编码器全连接；解码器因果；跨注意力 & 条件序列生成 & 翻译、摘要、结构化转换 \\
\bottomrule
\end{tabularx}
\caption{术语消化：模型家族由可见信息与训练目标共同定义。}
\end{table}

\teachervoice{Q\&A 中 Karpathy 说自己不太喜欢每次采样一个 token 就永久提交的自回归过程，更偏爱类似 diffusion 的“先草稿、再修订”。这是 2023 年的个人研究直觉，不是本讲证明的 Transformer 必然演化方向。}

\subsection{本章小结}

Decoder-only 用 causal self-attention 做 next-token generation；encoder-only 允许双向通信并采用不泄漏标签的表示学习目标；encoder--decoder 让 causal decoder 通过 cross-attention 查询源 memory。三者共享 attention 数学，只在 connectivity、state source 与 objective 上重新配置。

